% page 406 du fac-simile (image p-405 du scan)
% bloc 149.9 x 250.6 mm ; marge gauche 8.0 mm
\pgc{-12.700mm}{0.000mm}{
\l{}
\l{}
\l{}
\l{\cel{3}398.}
\l{}
\l{\sou{nuda}. I. Qua ne\cel{2}weras vesti. - II. Qua ne esas garnisita}
\l{\cel{3}(ye ulo). - EFIS.}
\l{}
\l{\sou{nudlo}. (\sou{koquarto}). Singla de la bendi streta qui konsistas}
\l{\cel{3}ye pasto ekfarino ed ovi, e tranchita por formacar ta}
\l{\cel{3}bendi. - DE.}
\l{}
\l{\sou{nugato}. Pasto ek mandeli torefaktita e sukro. - EF.}
\l{}
\l{\sou{nuko.}\cel{2}(\sou{anat.}) Parto dopa di kolo, ube ica sujuntas}
\l{\cel{3}a la kapo. - FIS.}
\l{}
\l{\sou{nukleo}. I. (\sou{biol.)}\cel{2}Parto kompakta ye qua konsistas la cen-}
\l{\cel{3}tro di ula kozi (protoplasmo, e c.). II. (\sou{Atomistiko}).}
\l{\cel{3}Korpo elektro-pozitiva,konsistanta ye elektro-pozitiva}
\l{\cel{4}protoni ed elektro-neutra neutroni, qua konstitucas la}
\l{\cel{3}maso precipua di atomi. - III. (\sou{metaf.}) Ta personi qui}
\l{\cel{3}su asemblis dat-unesme e cirkum qui plusi venas aglome-}
\l{\cel{3}rar su. - EFS.}
\l{}
\l{\sou{nukleino}. (\sou{biol.})\cel{2}Substanco ye qua konsistas la celulo-}
\l{\cel{3}nukleo e qua kontenas til 5/l00 fosforo. - DEFIS.}
\l{}
\l{\sou{nukleolo}. (\sou{biol.}) Mikra korpo ronda, unika o multopla, qua}
\l{\cel{3}existas en la nukleo di la celuli.}
\l{}
\l{\sou{nukleoplasmo}. (\sou{biol.)}\cel{2}Protoplasmo di la nukleo.}
\l{}
\l{\sou{nula}. Ne mem nur un.}
\l{}
\l{\sou{numero}.\cel{2}(\sou{aritm.})\cel{2}Vorto qua, en serio de termini, indikas}
\l{\cel{3}la rango di un ek la termini relate la ceteri. - DEFIRS.}
\l{}
\l{\sou{numeratoro}. (\sou{aritm.)}\cel{2}Ta ek la du nombri di fraciono, qua}
\l{\cel{3}indikas quanta partin inter-egala di la unajo lu konte-}
\l{\cel{3}nas. - DEFIRS.}
\l{}
\l{\sou{numismatiko}.\cel{2}La cienco qua traktas la moneto-peci e me-}
\l{\cel{3}dalii. DEFIRS.}
\l{}
\l{\sou{numulario}. (\sou{bot.)}\cel{1}- Planto, vicina ye la primrozo. L.\cel{1}\sou{Lysi}-}
\l{\cel{3}\sou{machia numularia}.\cel{2}- FIS.}
\l{}
\l{\sou{numulito}. (\sou{paleont.}) Genero de foraminiferi, tipo di la fa-}
\l{\cel{3}milio \textquotedbl{}numilitidei\textquotedbl{}, kun konki ronda, disko-forma, e di-}
\l{\cel{3}vidita ye mikra lojii multega. - DEFIS.}
\l{}
\l{\sou{nun}.\cel{2}En la instanto, en la epoko, en la periodo, prezenta.}
\l{\cel{3}- D. Gr.}
\l{}
\l{\sou{nuncio}.\cel{2}Ambasadisto di la papo. DEFIRS.}
\l{}
\l{\sou{nur}.\cel{2}Sen ulo od ulu plusa.\cel{2}D.}
}
